% Source: book pp. 12--17 (PDF 26--31); solutions pp. 9--13.
\section{Definition of Vector Space}

The motivation for the definition of a vector space comes from properties of
addition and scalar multiplication in $\F^n$: Addition is commutative,
associative, and has an identity. Every element has an additive inverse. Scalar
multiplication is associative. Scalar multiplication by $1$ acts as expected.
Addition and scalar multiplication are connected by distributive properties.

We will define a vector space to be a set $V$ with an addition and a scalar
multiplication on $V$ that satisfy the properties in the paragraph above.

\begin{definition}[Addition, scalar multiplication]\label{ax:1.19}\leavevmode
\begin{itemize}
\item An \emph{addition} on a set $V$ is a function that assigns an element
  $u+v\in V$ to each pair of elements $u,v\in V$.
\item A \emph{scalar multiplication} on a set $V$ is a function that assigns an
  element $\lambda v\in V$ to each $\lambda\in\F$ and each $v\in V$.
\end{itemize}
\end{definition}

Now we are ready to give the formal definition of a vector space.

\begin{definition}[Vector space]\label{ax:1.20}
A \emph{vector space} is a set $V$ along with an addition on $V$ and a scalar
multiplication on $V$ such that the following properties hold.
\begin{axprops}
\item[commutativity] $u+v=v+u$ for all $u,v\in V$.
\item[associativity] $(u+v)+w=u+(v+w)$ and $(ab)v=a(bv)$ for all
  $u,v,w\in V$ and for all $a,b\in\F$.
\item[additive identity] There exists an element $0\in V$ such that $v+0=v$
  for all $v\in V$.
\item[additive inverse] For every $v\in V$, there exists $w\in V$ such that
  $v+w=0$.
\item[multiplicative identity] $1v=v$ for all $v\in V$.
\item[distributive properties] $a(u+v)=au+av$ and $(a+b)v=av+bv$ for all
  $a,b\in\F$ and all $u,v\in V$.
\end{axprops}
\end{definition}

The following geometric language sometimes aids our intuition.

\begin{definition}[Vector, point]\label{ax:1.21}
Elements of a vector space are called \emph{vectors} or \emph{points}.
\end{definition}

The scalar multiplication in a vector space depends on $\F$. Thus when we need
to be precise, we will say that $V$ is a \emph{vector space over} $\F$ instead
of saying simply that $V$ is a vector space. For example, $\R^n$ is a vector
space over $\R$, and $\C^n$ is a vector space over $\C$.

\begin{definition}[Real vector space, complex vector space]\label{ax:1.22}\leavevmode
\begin{itemize}
\item A vector space over $\R$ is called a \emph{real vector space}.
\item A vector space over $\C$ is called a \emph{complex vector space}.
\end{itemize}
\end{definition}

Usually the choice of $\F$ is either clear from the context or irrelevant. Thus
we often assume that $\F$ is lurking in the background without specifically
mentioning it.

\marginnote{The simplest vector space is $\{0\}$, which contains only one
point.}%
With the usual operations of addition and scalar multiplication, $\F^n$ is a
vector space over $\F$, as you should verify. The example of $\F^n$ motivated
our definition of vector space.

\begin{example}[$\F^\infty$]\label{ax:1.23}
$\F^\infty$ is defined to be the set of all sequences of elements of $\F$:
\[ \F^\infty = \{(x_1,x_2,\dots) : x_k\in\F \text{ for } k=1,2,\dots\}. \]
Addition and scalar multiplication on $\F^\infty$ are defined as expected:
\begin{align*}
(x_1,x_2,\dots)+(y_1,y_2,\dots) &= (x_1+y_1,x_2+y_2,\dots),\\
\lambda(x_1,x_2,\dots) &= (\lambda x_1,\lambda x_2,\dots).
\end{align*}
With these definitions, $\F^\infty$ becomes a vector space over $\F$, as you
should verify. The additive identity in this vector space is the sequence of
all $0$'s.
\end{example}

Our next example of a vector space involves a set of functions.

\begin{notation}[$\F^S$]\label{ax:1.24}\leavevmode
\begin{itemize}
\item If $S$ is a set, then $\F^S$ denotes the set of functions from $S$ to
  $\F$.
\item For $f,g\in\F^S$, the \emph{sum} $f+g\in\F^S$ is the function defined by
  \[ (f+g)(x) = f(x)+g(x) \]
  for all $x\in S$.
\item For $\lambda\in\F$ and $f\in\F^S$, the \emph{product} $\lambda f\in\F^S$
  is the function defined by
  \[ (\lambda f)(x) = \lambda f(x) \]
  for all $x\in S$.
\end{itemize}
\end{notation}

As an example of the notation above, if $S$ is the interval $[0,1]$ and
$\F=\R$, then $\R^{[0,1]}$ is the set of real-valued functions on the interval
$[0,1]$.

You should verify all three bullet points in the next example.

\begin{example}[$\F^S$ is a vector space]\label{ax:1.25}\leavevmode
\begin{itemize}
\item If $S$ is a nonempty set, then $\F^S$ (with the operations of addition
  and scalar multiplication as defined above) is a vector space over $\F$.
\item The additive identity of $\F^S$ is the function $0\colon S\to\F$ defined
  by
  \[ 0(x) = 0 \]
  for all $x\in S$.
\item For $f\in\F^S$, the additive inverse of $f$ is the function
  $-f\colon S\to\F$ defined by
  \[ (-f)(x) = -f(x) \]
  for all $x\in S$.
\end{itemize}
\end{example}

\marginnote{The elements of the vector space $\R^{[0,1]}$ are real-valued
functions on $[0,1]$, not lists. In general, a vector space is an abstract
entity whose elements might be lists, functions, or weird objects.}%
{\emergencystretch=2em
The vector space $\F^n$ is a special case of the vector space $\F^S$ because
each $(x_1,\dots,x_n)\in\F^n$ can be thought of as a function $x$ from the set
$\{1,2,\dots,n\}$ to $\F$ by writing $x(k)$ instead of $x_k$ for the
$k^{\text{th}}$ coordinate of $(x_1,\dots,x_n)$. In other words, we can think
of $\F^n$ as $\F^{\{1,2,\dots,n\}}$. Similarly, we can think of $\F^\infty$ as
$\F^{\{1,2,\dots\}}$.\par}

Soon we will see further examples of vector spaces, but first we need to
develop some of the elementary properties of vector spaces.

The definition of a vector space requires it to have an additive identity. The
next result states that this identity is unique.

\begin{result}[Unique additive identity]\label{ax:1.26}
A vector space has a unique additive identity.
\end{result}

\begin{proof}
Suppose $0$ and $0'$ are both additive identities for some vector space $V$.
Then
\[ 0' = 0'+0 = 0+0' = 0, \]
where the first equality holds because $0$ is an additive identity, the second
equality comes from commutativity, and the third equality holds because $0'$
is an additive identity. Thus $0'=0$, proving that $V$ has only one additive
identity.
\end{proof}

Each element $v$ in a vector space has an additive inverse, an element $w$ in
the vector space such that $v+w=0$. The next result shows that each element in
a vector space has only one additive inverse.

\begin{result}[Unique additive inverse]\label{ax:1.27}
Every element in a vector space has a unique additive inverse.
\end{result}

\begin{proof}
Suppose $V$ is a vector space. Let $v\in V$. Suppose $w$ and $w'$ are additive
inverses of $v$. Then
\[ w = w+0 = w+(v+w') = (w+v)+w' = 0+w' = w'. \]
Thus $w=w'$, as desired.
\end{proof}

Because additive inverses are unique, the following notation now makes sense.

\begin{notation}[$-v$, $w-v$]\label{ax:1.28}
Let $v,w\in V$. Then
\begin{itemize}
\item $-v$ denotes the additive inverse of $v$;
\item $w-v$ is defined to be $w+(-v)$.
\end{itemize}
\end{notation}

Almost all results in this book involve some vector space. To avoid having to
restate frequently that $V$ is a vector space, we now make the necessary
declaration once and for all.

\begin{notation}[$V$]\label{ax:1.29}
For the rest of this book, $V$ denotes a vector space over $\F$.
\end{notation}

In the next result, $0$ denotes a scalar (the number $0\in\F$) on the left side
of the equation and a vector (the additive identity of $V$) on the right side
of the equation.

\begin{result}[The number $0$ times a vector]\label{ax:1.30}
$0v=0$ for every $v\in V$.
\end{result}

\begin{proof}
For $v\in V$, we have
\[ 0v = (0+0)v = 0v+0v. \]
Adding the additive inverse of $0v$ to both sides of the equation above gives
$0=0v$, as desired.
\end{proof}

\marginnote[-14mm]{The result in \ref{ax:1.30} involves the additive identity of $V$
and scalar multiplication. The only part of the definition of a vector space
that connects vector addition and scalar multiplication is the distributive
property. Thus the distributive property must be used in the proof of
\ref{ax:1.30}.}%
In the next result, $0$ denotes the additive identity of $V$. Although their
proofs are similar, \ref{ax:1.30} and \ref{ax:1.31} are not identical. More
precisely, \ref{ax:1.30} states that the product of the scalar $0$ and any
vector equals the vector $0$, whereas \ref{ax:1.31} states that the product of
any scalar and the vector $0$ equals the vector $0$.

\begin{result}[A number times the vector $0$]\label{ax:1.31}
$a0=0$ for every $a\in\F$.
\end{result}

\begin{proof}
For $a\in\F$, we have
\[ a0 = a(0+0) = a0+a0. \]
Adding the additive inverse of $a0$ to both sides of the equation above gives
$0=a0$, as desired.
\end{proof}

Now we show that if an element of $V$ is multiplied by the scalar $-1$, then
the result is the additive inverse of the element of $V$.

\begin{result}[The number $-1$ times a vector]\label{ax:1.32}
$(-1)v=-v$ for every $v\in V$.
\end{result}

\begin{proof}
For $v\in V$, we have
\[ v+(-1)v = 1v+(-1)v = \bigl(1+(-1)\bigr)v = 0v = 0. \]
This equation says that $(-1)v$, when added to $v$, gives $0$. Thus $(-1)v$ is
the additive inverse of $v$, as desired.
\end{proof}

% ------------------------------------------------------------------ exercises
\exercisesheading{1B}

\begin{exercise}
Prove that $-(-v)=v$ for every $v\in V$.
\end{exercise}
\begin{solution}
Let $v\in V$. By the definition of additive inverse, we have
\[ v+(-v) = 0 = -(-v)+(-v). \]
Since additive inverses are unique, we conclude that $-(-v)=v$.
\end{solution}

\begin{exercise}
Suppose $a\in\F$, $v\in V$, and $av=0$. Prove that $a=0$ or $v=0$.
\end{exercise}
\begin{solution}
If $a=0$, then we are done. Suppose $a\ne0$. Then
\[ v = 1v = \frac1a(av) = \frac1a\cdot0 = 0. \qedhere \]
\end{solution}

\begin{exercise}
Suppose $v,w\in V$. Explain why there exists a unique $x\in V$ such that
$v+3x=w$.
\end{exercise}
\begin{solution}
\solnote{Manual Remark}{The solution says ``dividing by $3$'' in the sense that
we're dividing by the scalar $3$ in the field $\F$, not a vector in $V$.}
Solving for $x$ explicitly, we obtain
\[ 3x = w-v \implies x = (1/3)(w-v). \]
This shows that there exists $x\in V$ such that $v+3x=w$. To see that $x$ is
unique, suppose there exists another $y\in V$ such that $v+3y=w$. Then dividing
by $3$, we get
\[ 3x = w-v = 3y \implies x = y. \qedhere \]
\end{solution}

\begin{exercise}
The empty set is not a vector space. The empty set fails to satisfy only one of
the requirements listed in the definition of a vector space (\ref{ax:1.20}).
Which one?
\end{exercise}
\begin{solution}
The empty set fails to satisfy the requirement that there exists a zero vector
$0\in V$ such that $v+0=v$ for all $v\in V$ since it contains no elements.
\end{solution}

\begin{exercise}
Show that in the definition of a vector space (\ref{ax:1.20}), the additive
inverse condition can be replaced with the condition that
\[ 0v=0 \text{ for all } v\in V. \]
Here the $0$ on the left side is the number $0$, and the $0$ on the right side
is the additive identity of $V$.
\exnote{The phrase a ``condition can be replaced'' in a definition means that
the collection of objects satisfying the definition is unchanged if the
original condition is replaced with the new condition.}
\end{exercise}
\begin{solution}
Suppose that $0v=0$ for all $v\in V$. Let $v\in V$. Then
\[ v+(-1)v = 1v+(-1)v = \bigl(1+(-1)\bigr)v = 0v = 0. \]
Hence, $(-1)v$ is the additive inverse of $v$.

Conversely, suppose that for every $v\in V$, there exists an additive inverse
$-v\in V$ such that $v+(-v)=0$. Then $0v=(0+0)v=0v+0v$. Adding $-0v$ to both
sides, we get $0=0v$. Hence, $0v=0$ for all $v\in V$.
\end{solution}

\begin{exercise}
Let $\infty$ and $-\infty$ denote two distinct objects, neither of which is in
$\R$. Define an addition and scalar multiplication on $\R\cup\{\infty,-\infty\}$
as you could guess from the notation. Specifically, the sum and product of two
real numbers is as usual, and for $t\in\R$ define
\[ t\infty = \begin{cases} -\infty & \text{if } t<0,\\ 0 & \text{if } t=0,\\
     \infty & \text{if } t>0, \end{cases}
   \qquad
   t(-\infty) = \begin{cases} \infty & \text{if } t<0,\\ 0 & \text{if } t=0,\\
     -\infty & \text{if } t>0, \end{cases} \]
and
\begin{align*}
t+\infty = \infty+t = \infty+\infty &= \infty,\\
t+(-\infty) = (-\infty)+t = (-\infty)+(-\infty) &= -\infty,\\
\infty+(-\infty) = (-\infty)+\infty &= 0.
\end{align*}
With these operations of addition and scalar multiplication, is
$\R\cup\{\infty,-\infty\}$ a vector space over $\R$? Explain.
\end{exercise}
\begin{solution}
No, $\R\cup\{\infty,-\infty\}$ is not a vector space over $\R$ as it does not
satisfy the associative property of addition. Let $t\in\R\setminus\{0\}$. Then
\[ (t+\infty)+(-\infty) = \infty+(-\infty) = 0 \ne t = t+0
   = t+\bigl(\infty+(-\infty)\bigr). \qedhere \]
\end{solution}

\begin{exercise}
Suppose $S$ is a nonempty set. Let $V^S$ denote the set of functions from $S$
to $V$. Define a natural addition and scalar multiplication on $V^S$, and show
that $V^S$ is a vector space with these definitions.
\end{exercise}
\begin{solution}
Let $f,g\in V^S$ and $t\in\F$. Define addition and scalar multiplication as
follows:
\[ (f+g)(s) = f(s)+g(s), \quad (tf)(s) = t\bigl(f(s)\bigr) \]
for all $s\in S$. We will show that $V^S$ is a vector space with these
operations.
\begin{enumerate}[label=\arabic*)]
\item \textbf{Commutativity}: Let $f,g\in V^S$. Then for all $s\in S$,
  \[ (f+g)(s) = f(s)+g(s) = g(s)+f(s) = (g+f)(s). \]
\item \textbf{Associativity}: Let $f,g,h\in V^S$ and $t,u\in\F$. Then for all
  $s\in S$,
  \begin{align*}
  \bigl((f+g)+h\bigr)(s) &= (f+g)(s)+h(s)\\
    &= \bigl(f(s)+g(s)\bigr)+h(s)\\
    &= f(s)+\bigl(g(s)+h(s)\bigr)\\
    &= f(s)+(g+h)(s)\\
    &= \bigl(f+(g+h)\bigr)(s),
  \end{align*}
  and
  \begin{align*}
  \bigl((tu)f\bigr)(s) &= (tu)\bigl(f(s)\bigr)\\
    &= t\bigl(u(f(s))\bigr)\\
    &= t(uf)(s)\\
    &= \bigl(t(uf)\bigr)(s).
  \end{align*}
\item \textbf{Additive identity}: Let $0\in V$ be the additive identity in $V$.
  Define the zero function $0_S\in V^S$ by $0_S(s)=0$ for all $s\in S$. Then
  for all $f\in V^S$ and $s\in S$,
  \[ (f+0_S)(s) = f(s)+0_S(s) = f(s)+0 = f(s) = 0+f(s) = (0_S+f)(s). \]
  Hence, $0_S$ is the additive identity in $V^S$.
\item \textbf{Additive inverse}: Let $f\in V^S$. Define the function
  $-f\in V^S$ by $(-f)(s)=-f(s)$ for all $s\in S$. Then for all $s\in S$,
  \[ \bigl(f+(-f)\bigr)(s) = f(s)+(-f)(s) = f(s)+\bigl(-f(s)\bigr) = 0
     = 0_S(s). \]
  Hence, $-f$ is the additive inverse of $f$ in $V^S$.
\item \textbf{Multiplicative identity}: For all $f\in V^S$ and $s\in S$,
  \[ (1f)(s) = 1\bigl(f(s)\bigr) = f(s). \]
\item \textbf{Distributive properties}: Let $f,g\in V^S$ and $t,u\in\F$. Then
  for all $s\in S$,
  \begin{align*}
  \bigl((t+u)f\bigr)(s) &= (t+u)\bigl(f(s)\bigr)\\
    &= t\bigl(f(s)\bigr)+u\bigl(f(s)\bigr)\\
    &= (tf)(s)+(uf)(s)\\
    &= (tf+uf)(s),
  \end{align*}
  and
  \begin{align*}
  \bigl(t(f+g)\bigr)(s) &= t\bigl((f+g)(s)\bigr)\\
    &= t\bigl(f(s)+g(s)\bigr)\\
    &= t\bigl(f(s)\bigr)+t\bigl(g(s)\bigr)\\
    &= (tf)(s)+(tg)(s)\\
    &= (tf+tg)(s). \qedhere
  \end{align*}
\end{enumerate}
\end{solution}

\begin{exercise}
Suppose $V$ is a real vector space.
\begin{itemize}
\item The \emph{complexification} of $V$, denoted by $V_{\C}$, equals
  $V\times V$. An element of $V_{\C}$ is an ordered pair $(u,v)$, where
  $u,v\in V$, but we write this as $u+iv$.
\item Addition on $V_{\C}$ is defined by
  \[ (u_1+iv_1)+(u_2+iv_2) = (u_1+u_2)+i(v_1+v_2) \]
  for all $u_1,v_1,u_2,v_2\in V$.
\item Complex scalar multiplication on $V_{\C}$ is defined by
  \[ (a+bi)(u+iv) = (au-bv)+i(av+bu) \]
  for all $a,b\in\R$ and all $u,v\in V$.
\end{itemize}
Prove that with the definitions of addition and scalar multiplication as
above, $V_{\C}$ is a complex vector space.
\exnote{Think of $V$ as a subset of $V_{\C}$ by identifying $u\in V$ with
$u+i0$. The construction of $V_{\C}$ from $V$ can then be thought of as
generalizing the construction of $\C^n$ from $\R^n$.}
\end{exercise}
\begin{solution}
We prove that $V_{\C}$ is a complex vector space by verifying the vector space
axioms.
\begin{enumerate}[label=\arabic*)]
\item \textbf{Commutativity}: Let $u_1+v_1i,u_2+v_2i\in V_{\C}$. Then
  \begin{align*}
  (u_1+v_1i)+(u_2+v_2i) &= (u_1+u_2)+i(v_1+v_2)\\
    &= (u_2+u_1)+i(v_2+v_1) = (u_2+v_2i)+(u_1+v_1i).
  \end{align*}
\item \textbf{Associativity}: Let $u_1+v_1i,u_2+v_2i,u_3+v_3i\in V_{\C}$. Then
  \begin{align*}
  \bigl((u_1+v_1i)+(u_2+v_2i)\bigr)+(u_3+v_3i)
    &= \bigl((u_1+u_2)+i(v_1+v_2)\bigr)+(u_3+v_3i)\\
    &= \bigl((u_1+u_2)+u_3\bigr)+i\bigl((v_1+v_2)+v_3\bigr)\\
    &= \bigl(u_1+(u_2+u_3)\bigr)+i\bigl(v_1+(v_2+v_3)\bigr)\\
    &= (u_1+v_1i)+\bigl((u_2+v_2i)+(u_3+v_3i)\bigr).
  \end{align*}
  Similarly, let $a+bi,c+di\in\C$ and $u+vi\in V_{\C}$. Then
  \begin{align*}
  &\bigl((a+bi)(c+di)\bigr)(u+vi)\\
    &\qquad= \bigl((ac-bd)+i(ad+bc)\bigr)(u+vi)\\
    &\qquad= \bigl((ac-bd)u-(ad+bc)v\bigr)+i\bigl((ac-bd)v+(ad+bc)u\bigr)\\
    &\qquad= \bigl(a(cu-dv)-b(du+cv)\bigr)+i\bigl(a(du+cv)+b(cu-dv)\bigr)\\
    &\qquad= (a+bi)\bigl((cu-dv)+i(du+cv)\bigr)\\
    &\qquad= (a+bi)\bigl((c+di)(u+vi)\bigr).
  \end{align*}
\item \textbf{Additive identity}: Let $0\in V$ be the additive identity in $V$.
  Then for all $u+vi\in V_{\C}$,
  \[ (u+vi)+(0+0i) = (u+0)+i(v+0) = u+vi = (0+0i)+(u+vi). \]
  Hence, $0+0i$ is the additive identity in $V_{\C}$.
\item \textbf{Additive inverse}: Let $u+vi\in V_{\C}$. Then
  \[ (u+vi)+\bigl(-u+(-v)i\bigr) = (u-u)+i(v-v) = 0+0i. \]
  Hence, $-u+(-v)i$ is the additive inverse of $u+vi$ in $V_{\C}$.
\item \textbf{Multiplicative identity}: For all $u+vi\in V_{\C}$,
  \[ (1+0i)(u+vi) = (1u-0v)+i(1v+0u) = u+vi. \]
\item \textbf{Distributive properties}: Let $a+bi,c+di\in\C$ and
  $u_1+v_1i,u_2+v_2i\in V_{\C}$. Then
  \begin{align*}
  &(a+bi)\bigl((u_1+v_1i)+(u_2+v_2i)\bigr)\\
    &\qquad= (a+bi)\bigl((u_1+u_2)+i(v_1+v_2)\bigr)\\
    &\qquad= \bigl(a(u_1+u_2)-b(v_1+v_2)\bigr)+i\bigl(a(v_1+v_2)+b(u_1+u_2)\bigr)\\
    &\qquad= (au_1-bv_1)+i(av_1+bu_1)+(au_2-bv_2)+i(av_2+bu_2)\\
    &\qquad= (a+bi)(u_1+v_1i)+(a+bi)(u_2+v_2i),
  \end{align*}
  and
  \begin{align*}
  &\bigl((a+bi)+(c+di)\bigr)(u_1+v_1i)\\
    &\qquad= \bigl((a+c)+(b+d)i\bigr)(u_1+v_1i)\\
    &\qquad= \bigl((a+c)u_1-(b+d)v_1\bigr)+i\bigl((a+c)v_1+(b+d)u_1\bigr)\\
    &\qquad= (au_1-bv_1)+i(av_1+bu_1)+(cu_1-dv_1)+i(cv_1+du_1)\\
    &\qquad= (a+bi)(u_1+v_1i)+(c+di)(u_1+v_1i). \qedhere
  \end{align*}
\end{enumerate}
\end{solution}
